\frontchapter{บทคัดย่อ}
เครื่องมือบริหารเส้นทางโจมตี (Attack Path Management) จัดอันดับเส้นทางที่อันตรายด้วยคะแนนที่อธิบายที่มาไม่ได้
และไม่วัดว่าอันดับเปลี่ยนเท่าไรเมื่อน้ำหนักคลาดเคลื่อน
โครงงานนี้สร้างกราฟการโจมตีของโครงสร้างพื้นฐานแบบ Hybrid ที่ทุกเส้นเชื่อม (edge) มีความน่าจะเป็น $p(e)$ อธิบายที่มาได้
(EPSS สำหรับ edge ที่มี CVE, ค่าที่ fit จากความสัมพันธ์ CVSS--EPSS สำหรับ misconfiguration
และ $p=1$ สำหรับความสัมพันธ์เชิงโครงสร้าง)
หาจุดตัดด้วย dominator tree เทียบกับ weighted node min-cut
และวัดความเสถียรของอันดับด้วย Jaccard@k และ Kendall's $\tau$

ทดลองบนกราฟจำลอง 216 โครงสร้าง (864 กราฟเมื่อรวมโมเดลน้ำหนัก 4 แบบ) และ lab ที่ใช้ Terraform, LocalStack, Cartography และ Neo4j
ผลคือ (1) อันดับทนต่อการรบกวนน้ำหนัก (Jaccard@5 เท่ากับ 0.954 ที่ $\pm 50\%$)
แต่ไม่ทนต่อ edge ที่ขาดหาย (0.49 เมื่อลบ 5\% และ 0.15 เมื่อลบ 20\%)
(2) จุดตัดที่การันตีทุกเส้นทางพบเพียง 18 จาก 216 กราฟ (8.3\%) ส่วนใหญ่ในเครือข่ายแบบ micro-segmentation
(3) weighted min-cut ที่ใช้น้ำหนัก EPSS อยู่ในเซตของ dominator tree เพียง 72\% ของกรณีที่เทียบได้
และฟังก์ชัน fit จาก CVSS อธิบายความแปรปรวนได้น้อย ($R^2$ ในสเกล logit เท่ากับ 0.219)

ข้อจำกัดคือกราฟจำลองไม่มีค่าจริงสำหรับเทียบความถูกต้อง จึงสรุปไม่ได้ว่า EPSS แม่นกว่า CVSS
และ lab ใช้ Docker, nmap และ LocalStack แทน GNS3, Nessus และ AWS จริง

\vspace{0.4cm}
\noindent\textbf{คำสำคัญ:} จุดตัด (Choke Point), กราฟการโจมตี, Dominator Tree, EPSS, CVSS, โครงสร้างพื้นฐานแบบ Hybrid
